\documentclass[11pt,a4paper]{article}
\usepackage[utf8]{inputenc}
\usepackage[T1]{fontenc}
\usepackage[portuguese]{babel}
\usepackage{geometry}
\geometry{margin=2.5cm}
\usepackage{listings}
\usepackage{xcolor}
\usepackage{amsmath}
\usepackage{amssymb}
\usepackage{hyperref}
\usepackage{enumitem}
\usepackage{tcolorbox}
\usepackage{tabularx}
\usepackage{booktabs}

\definecolor{codebg}{rgb}{0.95,0.95,0.95}
\definecolor{codegreen}{rgb}{0.1,0.5,0.1}
\definecolor{codepurple}{rgb}{0.5,0.1,0.5}
\definecolor{codegray}{rgb}{0.5,0.5,0.5}

\lstdefinestyle{python}{
    backgroundcolor=\color{codebg},
    commentstyle=\color{codegray}\itshape,
    keywordstyle=\color{codepurple}\bfseries,
    stringstyle=\color{codegreen},
    basicstyle=\ttfamily\small,
    breaklines=true,
    frame=single,
    language=Python,
    showstringspaces=false,
    tabsize=4,
    literate={á}{{\'a}}1 {ã}{{\~a}}1 {é}{{\'e}}1 {í}{{\'i}}1 {ó}{{\'o}}1 {õ}{{\~o}}1 {ú}{{\'u}}1 {ç}{{\c{c}}}1 {Á}{{\'A}}1 {Ã}{{\~A}}1 {É}{{\'E}}1 {Í}{{\'I}}1 {Ó}{{\'O}}1 {Õ}{{\~O}}1 {Ú}{{\'U}}1 {Ç}{{\c{C}}}1 {â}{{\^a}}1 {ê}{{\^e}}1 {ô}{{\^o}}1 {û}{{\^u}}1 {Â}{{\^A}}1 {Ê}{{\^E}}1 {Ô}{{\^O}}1 {Û}{{\^U}}1 {à}{{\`a}}1 {è}{{\`e}}1 {ì}{{\`i}}1 {ò}{{\`o}}1 {ù}{{\`u}}1 {ä}{{\"a}}1 {ö}{{\"o}}1 {Ä}{{\"A}}1 {Ö}{{\"O}}1
}
\lstset{style=python}

\newtcolorbox{objetivos}[1][]{
  colback=green!5!white,
  colframe=green!40!black,
  title=O que você vai aprender nesta página,
  fonttitle=\bfseries,
  #1
}

\newtcolorbox{exercicio}[1]{
  colback=blue!5!white,
  colframe=blue!40!black,
  title=#1,
  fonttitle=\bfseries
}

\newtcolorbox{solucao}{
  colback=yellow!5!white,
  colframe=orange!60!black,
  title=Solução proposta,
  fonttitle=\bfseries
}

\title{\textbf{Data Analysis with Python 2024--2025}\\Parte 7: Project Work\\\large Visão Geral e Instruções}
\author{Tradução e Adaptação: University of Helsinki --- MOOC.fi}
\date{}

\begin{document}

\maketitle

\section*{Nota Importante}
Este material é uma tradução e adaptação baseada no curso \textit{Data Analysis with Python 2024-2025}, oferecido pela \textbf{Universidade de Helsinque (MOOC.fi)}.

\tableofcontents
\newpage

\section{Introdução ao Projeto Final (Project Work)}

Nesta etapa final do curso, o estudante deve escolher e realizar um trabalho de projeto prático (\emph{project work}). Existem três opções principais de projetos disponíveis:
\begin{enumerate}
    \item \textbf{Análise de Sequências Biológicas} (\textit{Biological Sequence Analysis} - Parte 7);
    \item \textbf{Análise de Regressão em Dados Médicos} (\textit{Regression Analysis on Medical Data} - Parte 8);
    \item \textbf{Análise de Dados de Fósseis} (\textit{Fossil Data Analysis}).
\end{enumerate}

Para iniciar qualquer um dos projetos, o aluno deve baixar os exercícios e/ou cadernos de anotações (\emph{notebooks}) através do servidor TMC, tentar resolver o máximo de exercícios possíveis de forma independente e submeter as soluções para validação automática. O servidor TMC atua como uma ferramenta de auxílio ao progresso e validação prévia antes do processo de revisão por pares (\emph{peer review}).

\section{Fluxos de Trabalho Individuais}

Cada projeto possui um fluxo de trabalho específico:

\subsection{1. Análise de Sequências (Sequence Analysis)}
O estudante baixa a parte correspondente via TMC após completar a quantidade exigida da parte anterior. Dentro da pasta \texttt{src}, há o notebook \texttt{project\_notebook\_sequence\_analysis.ipynb}. O aluno deve preencher as células conforme as instruções e testar a integridade usando o comando \texttt{tmc test}. 

\begin{quote}
\textbf{Nota sobre os relatórios:} Ao lado de cada exercício no relatório, existem caixas de texto para preenchimento. A primeira serve para descrever, com suas próprias palavras, a ideia lógica da solução; a segunda serve para analisar os resultados obtidos com os dados de entrada.
\end{quote}

\subsection{2. Análise de Regressão (Regression Analysis)}
Consiste em ler a introdução teórica do documento específico de regressão, preenchendo diretamente as soluções nas células fornecidas no notebook sem alterar as linhas que contêm comentários de marcação de exercício (ex: \texttt{\# exercise x}), pois os testes automatizados dependem estritamente delas.

\subsection{3. Análise de Fósseis (Fossil Data Analysis)}
Siga rigorosamente as instruções do guia de atribuição específica do tema de fósseis fornecido nos materiais suplementares da plataforma MOOC.

\section{Processo de Revisão por Pares (Peer Review)}
Após a finalização e submissão do arquivo do projeto em formato Jupyter Notebook (\texttt{.ipynb}) na plataforma MOOC, o estudante precisará realizar a avaliação e o feedback de relatórios submetidos por outros três colegas.

\subsection{Executando Testes de Outros Alunos}
Se desejar testar o código de um colega durante a revisão por pares para verificar a exatidão das soluções:
\begin{itemize}
    \item Vá para uma pasta temporária (como \texttt{/tmp} em sistemas Unix) para evitar sobrescrever seu próprio trabalho.
    \item Baixe os testes usando o comando do cliente TMC adequado à versão da turma.
    \item Substitua o notebook do estudante na pasta de origem correspondente e execute o comando \texttt{tmc test}.
\end{itemize}

\subsection{Envio via TMC Paste}
Para enviar o projeto para a plataforma de avaliação:
\begin{itemize}
    \item \textbf{Via Cliente Gráfico (VS Code / TMC Client):} Clique com o botão direito no notebook do projeto e escolha \textit{Send Exercise to TMC Paste}. Copie o link gerado (iniciado com \url{https://tmc.mooc.fi/paste/}) e cole-o no formulário oficial da atividade.
    \item \textbf{Via CLI (Terminal):} Utilize o comando \texttt{tmc paste} dentro da pasta do exercício correspondente.
\end{itemize}

\end{document}
